\chapter[Теорема о неявной функции, заданной одним уравнением.]{Теорема о неявной функции, заданной одним уравнением.}

\section{Теорема о неявной функции, заданной одним уравнением}
\begin{defn}
Функция $y=f(x)$ называется \textit{неявной функцией, заданной уравнением}\rindex{функция!неявная, заданная одним уравнением} $F(x,y)=0$, если $F(x,f(x))=0$ для любого $x\in D_f$. (Здесь, как обычно, $D_f$ "--- область определения функции $f$.)
\end{defn}

%----------------------------------------------------------------------------------------
% Теорема 1 --- О неявной функции
%----------------------------------------------------------------------------------------

\begin{thm} \label{th1}
Пусть функция $F$ двух переменных $x$ и $y$ удовлетворяет следующим условиям:
\begin{enumerate}
    \item[$1^\circ$] $F$ непрерывна на некоторой окрестности $U(x_0, y_0)$ точки $(x_0, y_0)$;
    \item[$2^\circ$] $F(x_0, y_0) = 0$;
    \item[$3^\circ$] $F'_y(x_0, y_0) \ne 0$, $F'_y$ непрерывна в точке $(x_0, y_0)$.
\end{enumerate}
Тогда существует прямоугольная окрестность $Q_{\delta, \varepsilon}(x_0, y_0) = Q_\delta(x_0) \times Q_\varepsilon(y_0)$ точки $(x_0, y_0)$ такая, что на ней
$$
F(x, y) = 0 \iff y = f(x),
$$
где функция
$$
f \colon Q_\delta(x_0) \to Q_\varepsilon(y_0)
$$
непрерывна на $Q_\delta(x_0)$, $f(x_0) = y_0$.

Если дополнительно считать, что
\begin{enumerate}
    \item[$4^\circ$] $F$ дифференцируема в точке $(x_0, y_0)$,
\end{enumerate}
то $f$ дифференцируема в точке $x_0$ и
$$
f'(x_0) = -\frac{F'_x(x_0, y_0)}{F'_y(x_0, y_0)}.
$$
Если же при этом частные производные $F'_x$, $F'_y$ непрерывны на $U(x_0, y_0)$, то производная $f'(x) = -\frac{F'_x(x, f(x))}{F'_y(x, f(x))}$ непрерывна на $Q_\delta(x_0)$.
\end{thm}

\begin{proof}
На самом деле утверждения теоремы верны, если её условие $3^\circ$ заменить более общим:
\begin{enumerate}
    \item[$3^{\circ\circ}$] $F'_y$ существует и сохраняет знак на $U(x_0, y_0)$.
\end{enumerate}
Мы будем проводить доказательство при условии $3^{\circ\circ}$ вместо $3^\circ$.

Пусть $\sigma, \varepsilon > 0$ настолько малы, что в замыкании прямоугольной окрестности $Q_{\sigma, \varepsilon}(x_0, y_0)$ функция $F$ непрерывна, а $F'_y$ сохраняет знак. Ради определённости будем считать, что $F'_y > 0$ на $Q_{\sigma, \varepsilon}(x_0, y_0)$. Поэтому при каждом фиксированном $x \in Q_\sigma(x_0)$ функция $F(x, y)$, рассматриваемая как функция переменного $y$, строго возрастает на отрезке $[y_0 - \varepsilon, y_0 + \varepsilon]$.

Отсюда следует (поскольку $F(x_0, y_0) = 0$), что
$$
F(x_0, y_0 - \varepsilon) < 0, \quad F(x_0, y_0 + \varepsilon) > 0.
$$

Функции $F(x, y_0 - \varepsilon)$, $F(x, y_0 + \varepsilon)$ как функции переменного $x$ непрерывны на $Q_\sigma(x_0)$ (и, следовательно, обладают свойством сохранения знака), так что найдётся $\delta \in (0, \sigma]$ такое, что
$$
F(x, y_0 - \varepsilon) < 0, \quad F(x, y_0 + \varepsilon) > 0 \quad \forall\, x \in Q_\delta(x_0).
$$

Зафиксируем произвольное значение $x^* \in Q_\delta(x_0)$. Поскольку $F(x^*, y_0 - \varepsilon) < 0$, $F(x^*, y_0 + \varepsilon) > 0$, то по теореме Коши о промежуточном значении непрерывной функции $F(x^*, y)$ найдётся $y^* \in (y_0 - \varepsilon, y_0 + \varepsilon)$, при котором $F(x^*, y^*) = 0$. Такое значение $y^*$ единственно в силу строгой монотонности функции $F(x^*, y)$.

\begin{center}
\textit{Рис. 2}
\end{center}

Обозначим $y^* = f(x^*)$. Таким образом, построена (рис. 2) функция $f \colon Q_\delta(x_0) \to Q_\varepsilon(y_0)$ такая, что $f(x_0) = y_0$,
$$
F(x, y) = 0 \iff y = f(x) \quad \text{на} \quad Q_{\delta, \varepsilon}(x_0, y_0).
$$

Из последнего соотношения получаем, что
$$
F(x, f(x)) = 0 \quad \text{при} \quad x \in Q_\delta(x_0).
$$

Установим непрерывность функции $f$ на $Q_\delta(x_0)$. Непрерывность $f$ в точке $x_0$ следует из того, что в приведённых построениях число $\varepsilon > 0$ можно считать сколь угодно малым, причём для каждого достаточно малого $\varepsilon > 0$ было указано $\delta = \delta(\varepsilon) > 0$ такое, что $f(Q_\delta(x_0)) \subset Q_\varepsilon(y_0) = Q_\varepsilon(f(x_0))$.

Пусть теперь $x^*$ "--- произвольная точка из $Q_\delta(x_0)$, $y^* = f(x^*)$. Очевидно, что если условия теоремы 1 с заменой $3^\circ$ на $3^{\circ\circ}$ выполняются, то они выполняются и после замены в них $(x_0, y_0)$ на $(x^*, y^*)$. Следовательно, по уже доказанному, $f$ непрерывна в точке $x^*$.

Предположим теперь, что выполняется условие $4^\circ$. В силу дифференцируемости функции $F$ и замечания 11.1.1
\begin{multline} \label{eq3}
F(x_0 + \Delta x, y_0 + \Delta y) - F(x_0, y_0) = F'_x(x_0, y_0)\Delta x + {} \\
{} + F'_y(x_0, y_0)\Delta y + \varepsilon_1(\Delta x, \Delta y)\Delta x + \varepsilon_2(\Delta x, \Delta y)\Delta y, \quad (3)
\end{multline}
где $\varepsilon_i(\Delta x, \Delta y) \to 0$ при $(\Delta x, \Delta y) \to (0, 0)$, $i = 1, 2$. Здесь будем считать $|\Delta x|$ достаточно малым, $\Delta y = f(x_0 + \Delta x) - f(x_0) = \Delta f(x_0)$, так что $y_0 + \Delta y = f(x_0 + \Delta x)$.

Тогда из (3) получаем
\begin{multline*}
0 = F'_x(x_0, y_0)\Delta x + F'_y(x_0, y_0)\Delta y + {} \\
{} + \varepsilon_1(\Delta x, \Delta y)\Delta x + \varepsilon_2(\Delta x, \Delta y)\Delta y.
\end{multline*}

Так как $|\Delta x|$, а значит, и $|\Delta y| = |\Delta f|$ достаточно малы, имеем
$$
\Delta y = - \frac{F'_x(x_0, y_0) + \varepsilon_1(\Delta x, \Delta y)}{F'_y(x_0, y_0) + \varepsilon_2(\Delta x, \Delta y)}\Delta x = - \frac{F'_x(x_0, y_0)}{F'_y(x_0, y_0)}\Delta x + o(\Delta x)
$$
при $\Delta x \to 0$. Следовательно, функция $f$ дифференцируема в точке $x_0$ и
$$
f'(x_0) = - \frac{F'_x(x_0, y_0)}{F'_y(x_0, y_0)} = - \frac{F'_x(x_0, f(x_0))}{F'_y(x_0, f(x_0))}.
$$

Если функция $F$ дифференцируема не только в точке $(x_0, y_0)$, но и на окрестности $Q_{\delta, \varepsilon}(x_0, y_0)$, то последняя формула верна для любого $x \in Q_\delta(x_0)$:
$$
f'(x) = - \frac{F'_x(x, f(x))}{F'_y(x, f(x))} \quad \forall\, x \in Q_\delta(x_0).
$$

Из этой формулы вытекает последнее утверждение теоремы.
\end{proof}

\begin{notion}
В формулировке утверждения теоремы 1 можно отказаться от требования
$$
f \colon Q_\delta(x_0) \to Q_\varepsilon(y_0).
$$
Тогда, уменьшив при необходимости $\varepsilon$ или $\delta$, можно взять $\varepsilon = \delta$. При этом сохраняется свойство эквивалентности
$$
F(x, y) = 0 \iff y = f(x)
$$
на $Q_\delta(x_0, y_0) = Q_\delta(x_0) \times Q_\delta(y_0) \subset \bbR^2$, тождество $F(x, f(x)) = 0$ при $x \in Q_\delta(x_0)$ и дополнительное утверждение теоремы 1.

Изменённая таким образом теорема 1 равносильна, очевидно, исходной теореме 1.
\end{notion}